\begin{theorem}[High stationary-class norm product]
\label{mod:parameters-highproduct}
Let $K/F$ be a totally ramified prime-cyclic extension with residue degree
one, and retain the lower-break, integral-uniformizer, and monogenicity
hypotheses of the high-parameter modules.  Thus the proofs retain
$f(K/F)=1$ and $e(K/F)=[K:F]$ in the norm and valuation directions.
Suppose the actual stationary conductor decompositions are
\[
 m_K=2d_K+\varepsilon_K,
 \qquad m_F=2d+\varepsilon,
 \qquad \varepsilon_K,\varepsilon\in\{0,1\},
 \qquad m_K,m_F>1,
\]
and let
\[
 h:\operatorname{NormPolynomialPrecision}
   (m_K,d_K,\varepsilon_K;m_F,d,\varepsilon)
\]
be the precision certificate supplied by the relevant high-conductor
branch.  Its ambiguity component is exactly the inclusion
\[
 N_{K/F}(U_K^{d_K})\subseteq U_F^d.
\]
Consequently it induces the honest multiplicative quotient map
\begin{equation}
\label{eq:lean-highproduct-unit-norm}
 \mathcal N_h:U_K^0/U_K^{d_K}\longrightarrow U_F^0/U_F^d,
 \qquad [u]\longmapsto[N_{K/F}(u)].
\end{equation}
This is the only norm of a stationary class constructed here.  In
particular, Lean does not define a norm on an arbitrary element or
representative of the additive quotients
\[
 \mathcal O_K/\mathfrak p_K^{d_K}
 \quad\hbox{and}\quad
 \mathcal O_F/\mathfrak p_F^d.
\]

The structure \texttt{HighStationaryNormProduct} consists of a supplied
upstairs unit $u_K\in U_K^0$, a finite family of supplied downstairs units
$u_i\in U_F^0$, and the quotient equality
\begin{equation}
\label{eq:lean-highproduct-quotient-product}
 \boxed{\quad
   \mathcal N_h([u_K])=\prod_i[u_i]
   \quad\text{in }U_F^0/U_F^d.
 \quad}
\end{equation}
Its field-valued reading is only the congruence
\[
 N_{K/F}(u_K)\equiv\prod_i u_i\pmod{\mathfrak p_F^d},
\]
not an equality in $F$.  An exact depth-$d$ congruence constructs this
structure through \texttt{highStationaryNormProduct\_of\_congruence}.
If $u_K$ is replaced by a unit congruent to it modulo
$\mathfrak p_K^{d_K}$, the left side of
\eqref{eq:lean-highproduct-quotient-product} is unchanged.
Independently replacing any $u_i$ by a unit congruent modulo
$\mathfrak p_F^d$ leaves the right side unchanged.  These are precisely
depth-$d_K$ and depth-$d$ assertions; no finer representative-independence
statement is made.

To connect \eqref{eq:lean-highproduct-quotient-product} with stationary
coefficients, Lean starts with
separate, supplied additive representatives and quotient certificates
\[
 [u_K]_{\mathrm{add}}
  =\operatorname{St}_{K}(\chi_K,\psi_K;\gamma_K)
   \quad\text{in }\mathcal O_K/\mathfrak p_K^{d_K},
\]
and
\[
 [u_i]_{\mathrm{add}}
  =\operatorname{St}_{F}(\chi_i,\psi_F;\gamma_i)
   \quad\text{in }\mathcal O_F/\mathfrak p_F^d.
\]
The stationary numerator theorem proves that every such supplied
representative has order zero, after which it may be wrapped as an element
of $U^0$.  This extracts a property of a supplied representative; it does
not select a distinguished lift of either additive quotient.

In the odd wild non-stable branch put
\[
 p=[K:F],\qquad T=t+1,
 \qquad T\le m_F<2T,
\]
and assume that $\chi_F$ is minimal in its complete norm-character orbit.
The completed intermediate theorem supplies independent subcritical norm
choices $\alpha_1,\beta_1\in K^\times$ only at the permitted precisions
$\lfloor T/2\rfloor$ and $d$.  Write
\[
 \alpha_N=N_{K/F}(\alpha_1),
 \qquad \beta_N=N_{K/F}(\beta_1).
\]
The package also retains the supplied downstairs coefficients $\alpha$ and
$\beta$ with exactly the certificates
\[
 \alpha-\alpha_N\in
   \mathfrak p_F^{\,m_F-T+\lfloor T/2\rfloor},
 \qquad
 \beta-\beta_N\in\mathfrak p_F^d;
\]
neither certificate is strengthened to an equality in $F$.
For the Teichmuller scalar $[j]$ attached to
$j\in\mathbf F_p$, the package \texttt{OddIntermediateHighProductData}
uses the supplied representatives
\[
 u_K=\iota_{F,K}(\beta_N)
      -\beta_1\frac{\iota_{F,K}(\alpha_N)}{\alpha_1},
 \qquad
 u_j=N_{K/F}(\beta_1+[j]\alpha_1).
\]
They satisfy the actual additive stationary-class equalities
\[
 [u_K]_{\mathrm{add}}
  =\operatorname{St}_{\gamma_K,d_K+\varepsilon_K}(\chi_K)
  \quad\text{in }\mathcal O_K/\mathfrak p_K^{d_K},
\]
and, for every $j$ including $j=0$,
\[
 [u_j]_{\mathrm{add}}
  =\operatorname{St}_{\gamma_F,d+\varepsilon}(\tau^j\chi_F)
  \quad\text{in }\mathcal O_F/\mathfrak p_F^d.
\]
The twist datum, its actual conductor decomposition, and the admissibility
of $\gamma_F$ for that actual conductor are retained for every index.  The
intermediate theorem proves, at exactly the target depth,
\[
 u_j\equiv\beta_N+[j]\alpha_N\pmod{\mathfrak p_F^d}.
\]
Thus the stored factor is the exact norm representative, while its linear
expression is only a congruent representative of the same quotient class.

Lean proves the Teichmuller product identity for
these supplied scalars and the symmetric-function bound for the normalized
upstairs representative.  The resulting closed congruence is
\[
 N_{K/F}(u_K)
  \equiv \beta_N^p-\beta_N\alpha_N^{p-1}
  =\prod_{j\in\mathbf F_p}(\beta_N+[j]\alpha_N)
  \pmod{\mathfrak p_F^d}
\]
and the factorwise congruences identify the last product with
$\prod_j u_j$ in the same depth-$d$ quotient.  Consequently,
\[
 \mathcal N_h([u_K])
   =\prod_{j\in\mathbf F_p}[u_j]
   \quad\text{in }U_F^0/U_F^d.
\]
The norm points from $K$ to $F$, while the stationary adjoint used to
certify $u_K$ points from the downstairs coefficient quotient to the
upstairs one and retains the ordered common-denominator pair
$(\gamma_F,\iota_{F,K}\gamma_F)$.

In the positive-break wild quadratic branch,
\texttt{WildQuadraticHighProductData} reuses the simultaneous package of
the quadratic theorem without changing its choices.  Here
$\beta_1$ is chosen only through the break-level decomposition below, and
$\alpha_1$ only at precision $\lfloor(t+1)/2\rfloor$.  Thus
\[
 \beta=\delta_hN_{K/F}(\beta_1),
 \qquad \delta_h\in U_F^t,
 \qquad \alpha_N=N_{K/F}(\alpha_1),
\]
and the essential corrected coefficient is
\[
 c_F=\delta_h\alpha_N.
\]
For its supplied upstairs and two downstairs representatives,
\[
 u_K=\beta-\frac{\beta_1c_F}{\alpha_1},
 \qquad u_0=\beta,
 \qquad u_1=\beta+c_F,
\]
Lean indexes the factors by \texttt{Bool}, with \texttt{true} storing $u_0$
and \texttt{false} storing $u_1$; its certified ordered product is
$u_0u_1=\beta(\beta+c_F)$.  Lean retains the actual additive quotient
equalities for $\chi_K$,
$\chi_F$, and $\tau\chi_F$.  The exact quadratic norm calculation gives
\[
 N_{K/F}(u_K)\equiv
   \beta(\beta+c_F)
   =u_0u_1\pmod{\mathfrak p_F^d},
\]
hence
\[
 \mathcal N_h([u_K])=[u_0][u_1]
   \quad\text{in }U_F^0/U_F^d.
\]
The factor $\delta_h$ is part of $c_F$ throughout and is never replaced by
$1$ or discarded.  As in the odd branch, this is a quotient equality for
simultaneously supplied representatives, together with independence at
exactly the certified ambiguity depths; it asserts neither a canonical
stationary representative nor a field equality.

The proposition \texttt{HighProductBranch} retains the two formula-different
interfaces as separate universally quantified fields: the odd field returns
the existence of \texttt{OddIntermediateHighProductData}, and the quadratic
field returns the existence of \texttt{WildQuadraticHighProductData}.  The
principal theorem
\texttt{highParameter\_product} fills those fields with
\texttt{highParameter\_product\_intermediate} and
\texttt{highParameter\_product\_wildQuadratic}.  It assumes neither the
public high table, low-parameter table, phase-reduction result, nor any case
theorem.

\LeanFile{LanglandsFirstMainLemma/Parameters/HighProduct.lean}
\lean{LanglandsFirstMainLemma.HighStationaryNormProduct}
\lean{LanglandsFirstMainLemma.HighStationaryNormProduct.congruentAtTarget}
\lean{LanglandsFirstMainLemma.HighStationaryNormProduct.upstairs_representative_independent}
\lean{LanglandsFirstMainLemma.HighStationaryNormProduct.factors_representative_independent}
\lean{LanglandsFirstMainLemma.highStationaryNormProduct_of_congruence}
\lean{LanglandsFirstMainLemma.stationaryCoefficientClass_suppliedRepresentative_ord}
\lean{LanglandsFirstMainLemma.OddIntermediateHighProductData}
\lean{LanglandsFirstMainLemma.highParameter_intermediate_norm_product_congruent}
\lean{LanglandsFirstMainLemma.highParameter_product_intermediate}
\lean{LanglandsFirstMainLemma.WildQuadraticHighProductData}
\lean{LanglandsFirstMainLemma.highParameter_product_wildQuadratic}
\lean{LanglandsFirstMainLemma.HighProductBranch}
\lean{LanglandsFirstMainLemma.highParameter_product}
\leanok
\SourceText{Proposition prop:high-parameter-table, equation eq:high-product-congruence, and Proposition prop:wild-quadratic-high-parameters, equation eq:high-product-two.}
\uses{mod:parameters-highintermediate,mod:parameters-highquadratic,mod:parameters-normpolynomial}
\FormalizationContract
The source and target additive stationary coefficient quotients are exactly
$\mathcal O_K/\mathfrak p_K^{d_K}$ and
$\mathcal O_F/\mathfrak p_F^d$, whereas norm is induced only on
$U_K^0/U_K^{d_K}\to U_F^0/U_F^d$ from the ambiguity component of the
actual \texttt{NormPolynomialPrecision} certificate.  The odd branch uses
the complete $\mathbf F_p$-indexed product in the manuscript's direction;
the quadratic branch keeps the correction unit $\delta_h\in U_F^t$ in
both its twist factor and its norm product.  Every stationary equality and
norm-product assertion is a quotient equality or the corresponding exact
depth congruence.  Independence is proved only at depths $d_K$ and $d$.
No additive quotient is arbitrarily normed, no quotient congruence is
strengthened to an equality in a field, and no canonical stationary
representative is defined.
\end{theorem}
